% GENERATED FILE. Edit canonical lessons, then run npm run generate:books.
% canonical-source-hash: fnv1a64:47d5034d96279bfe
% canonical-lessons: IT-C339-opinione, IT-C339-parere, IT-C339-convinto, IT-C339-sostenere, IT-C339-punto-di-vista

\chapter{Secondo me}
\label{ch:it-secondo-me}
\hlchaptermodality{339}

\begin{chapteropening}
\textbf{By the end of this chapter:} \emph{I can give my opinion and my point of view, say that I am convinced and that I support an idea.}

Complete the last lesson of chapter 339: I can give my opinion and my point of view, say that I am convinced and that I support an idea.
\end{chapteropening}

\section[l'opinione]{l'opinione — the opinion}

\label{lesson:IT-C339-opinione}



\begin{quote}
Before the new one: say the Italian for exactly, then the Italian for almost.
\end{quote}

\subsection*{You'll want to know: l'opinione}
\textbf{l'opinione} — \textquotedblleft{}the opinion\textquotedblright{}.

\begin{quote}
Ognuno ha la sua opinione.
\end{quote}

\begin{quote}
\emph{Everyone has their own opinion.}
\end{quote}

\emph{Secondo me} opens an opinion; \emph{cambiare opinione} is to change your mind.

\begin{cousinweb}
\emph{opinione} is a learned borrowing from Latin \emph{opīniō, opīniōnis} \textquotedblleft{}belief, supposition, opinion\textquotedblright{} (Lewis \& Short), built on the verb \emph{opīnārī} \textquotedblleft{}to suppose\textquotedblright{}.
\end{cousinweb}

\begin{grammarlens}[title={What you've built}]
One more word for telling, explaining and describing at B1.
\end{grammarlens}

\subsection*{Your turn}
\begin{itemize}
\raggedright
  \item \emph{Say it:} \emph{l'opinione}
  \item \emph{Say it:} \emph{l'opinione}, once more
  \item \emph{Say it:} a sentence of your own with \emph{l'opinione}
  \item \emph{Recall:} say the Italian for exactly, then the Italian for almost, then say \emph{l'opinione} again
\end{itemize}

\begin{tcolorbox}[breakable,colback=teal!4,colframe=teal!35!black,title={Before you move on}]
What does l'opinione mean? (\textquotedblleft{}The opinion\textquotedblright{}.) Say it once more.
\end{tcolorbox}
\section[il parere]{il parere — the view (someone holds), the considered opinion}

\label{lesson:IT-C339-parere}



\begin{quote}
Before the new one: say the Italian for almost, then the Italian for the opinion.
\end{quote}

\subsection*{You'll want to know: il parere}
\textbf{il parere} — \textquotedblleft{}the view (someone holds), the considered opinion\textquotedblright{}.

\begin{quote}
Vorrei sentire il tuo parere prima di decidere.
\end{quote}

\begin{quote}
\emph{I would like to hear your view before deciding.}
\end{quote}

Also the verb \emph{parere} \textquotedblleft{}to seem\textquotedblright{}: \emph{mi pare giusto} — it seems right to me. \emph{A mio parere} is a little more formal than \emph{secondo me}.

\begin{cousinweb}
\emph{il parere} is the verb \emph{parere} \textquotedblleft{}to seem\textquotedblright{} used as a noun: what seems right to someone. The verb is inherited from Latin \emph{pārēre} \textquotedblleft{}to appear, to come into view\textquotedblright{} (Lewis \& Short).
\end{cousinweb}

\begin{grammarlens}[title={What you've built}]
One more word for telling, explaining and describing at B1.
\end{grammarlens}

\subsection*{Your turn}
\begin{itemize}
\raggedright
  \item \emph{Say it:} \emph{il parere}
  \item \emph{Say it:} \emph{il parere}, once more
  \item \emph{Say it:} a sentence of your own with \emph{il parere}
  \item \emph{Recall:} say the Italian for almost, then the Italian for the opinion, then say \emph{il parere} again
\end{itemize}

\begin{tcolorbox}[breakable,colback=teal!4,colframe=teal!35!black,title={Before you move on}]
What does il parere mean? (\textquotedblleft{}The view (someone holds), the considered opinion\textquotedblright{}.) Say it once more.
\end{tcolorbox}
\section[convinto]{convinto — convinced, sure}

\label{lesson:IT-C339-convinto}



\begin{quote}
Before the new one: say the Italian for the opinion, then the Italian for the view.
\end{quote}

\subsection*{You'll want to know: convinto}
\textbf{convinto} — \textquotedblleft{}convinced, sure\textquotedblright{}.

\begin{quote}
Sono convinto che sia la scelta giusta.
\end{quote}

\begin{quote}
\emph{I am convinced it is the right choice.}
\end{quote}

After \emph{sono convinto che} Italian uses the subjunctive: \emph{che sia}. A woman says \emph{convinta}.

\begin{cousinweb}
\emph{convinto} is the past participle of \emph{convincere} \textquotedblleft{}to convince\textquotedblright{}, which comes from Latin \emph{convincere} \textquotedblleft{}to overcome in argument, to prove\textquotedblright{} (Lewis \& Short), built on \emph{vincere} \textquotedblleft{}to conquer\textquotedblright{}.
\end{cousinweb}

\begin{grammarlens}[title={What you've built}]
One more word for telling, explaining and describing at B1.
\end{grammarlens}

\subsection*{Your turn}
\begin{itemize}
\raggedright
  \item \emph{Say it:} \emph{convinto}
  \item \emph{Say it:} \emph{convinto}, once more
  \item \emph{Say it:} a sentence of your own with \emph{convinto}
  \item \emph{Recall:} say the Italian for the opinion, then the Italian for the view, then say \emph{convinto} again
\end{itemize}

\begin{tcolorbox}[breakable,colback=teal!4,colframe=teal!35!black,title={Before you move on}]
What does convinto mean? (\textquotedblleft{}Convinced, sure\textquotedblright{}.) Say it once more.
\end{tcolorbox}
\section[sostenere]{sostenere — to maintain, to claim; to support}

\label{lesson:IT-C339-sostenere}



\begin{quote}
Before the new one: say the Italian for the view, then the Italian for convinced.
\end{quote}

\subsection*{You'll want to know: sostenere}
\textbf{sostenere} — \textquotedblleft{}to maintain, to claim; to support\textquotedblright{}.

\begin{quote}
Lui sostiene che il treno è più comodo dell'aereo.
\end{quote}

\begin{quote}
\emph{He maintains that the train is more comfortable than the plane.}
\end{quote}

Conjugated like \emph{tenere}: \emph{sostengo, sostieni, sostiene}. It can also mean to hold up or support a person or a weight.

\begin{cousinweb}
\emph{sostenere} is inherited from Latin \emph{sustinēre} \textquotedblleft{}to hold up, to bear\textquotedblright{} (Lewis \& Short), built from \emph{sub} \textquotedblleft{}under\textquotedblright{} and \emph{tenēre} \textquotedblleft{}to hold\textquotedblright{}: you hold an idea up so it does not fall.
\end{cousinweb}

\begin{grammarlens}[title={What you've built}]
One more word for telling, explaining and describing at B1.
\end{grammarlens}

\subsection*{Your turn}
\begin{itemize}
\raggedright
  \item \emph{Say it:} \emph{sostenere}
  \item \emph{Say it:} \emph{sostenere}, once more
  \item \emph{Say it:} a sentence of your own with \emph{sostenere}
  \item \emph{Recall:} say the Italian for the view, then the Italian for convinced, then say \emph{sostenere} again
\end{itemize}

\begin{tcolorbox}[breakable,colback=teal!4,colframe=teal!35!black,title={Before you move on}]
What does sostenere mean? (\textquotedblleft{}To maintain, to claim; to support\textquotedblright{}.) Say it once more.
\end{tcolorbox}
\section[il punto di vista]{il punto di vista — the point of view}

\label{lesson:IT-C339-punto-di-vista}



\begin{quote}
Before the new one: say the Italian for convinced, then the Italian for to maintain.
\end{quote}

\subsection*{You'll want to know: il punto di vista}
\textbf{il punto di vista} — \textquotedblleft{}the point of view\textquotedblright{}.

\begin{quote}
Dal mio punto di vista, è una buona idea.
\end{quote}

\begin{quote}
\emph{From my point of view, it is a good idea.}
\end{quote}

\emph{Dal mio punto di vista} — from where I stand. The plural is \emph{i punti di vista}.

\begin{cousinweb}
\emph{punto} is inherited from Latin \emph{punctum} \textquotedblleft{}a prick, a point\textquotedblright{} (Lewis \& Short), and \emph{vista} is built on \emph{vedere}, inherited from Latin \emph{vidēre} \textquotedblleft{}to see\textquotedblright{}: the spot you look from.
\end{cousinweb}

\begin{grammarlens}[title={What you've built}]
One more word for telling, explaining and describing at B1.
\end{grammarlens}

\subsection*{Your turn}
\begin{itemize}
\raggedright
  \item \emph{Say it:} \emph{il punto di vista}
  \item \emph{Say it:} \emph{il punto di vista}, once more
  \item \emph{Say it:} a sentence of your own with \emph{il punto di vista}
  \item \emph{Recall:} say the Italian for convinced, then the Italian for to maintain, then say \emph{il punto di vista} again
\end{itemize}

\begin{tcolorbox}[breakable,colback=teal!4,colframe=teal!35!black,title={Before you move on}]
What does il punto di vista mean? (\textquotedblleft{}The point of view\textquotedblright{}.) Say it once more.
\end{tcolorbox}
